% ==============================================================================
% Appendix: SFE Derivations
% ==============================================================================

\section{Supply Function Equilibrium Derivations}
\label{app:derivations}

This appendix derives the supply function equilibrium first-order
condition~\eqref{eq:sfe_foc} stated in Section~\ref{sec:model}, the
symmetric reduction~\eqref{eq:sfe_ode}, and the Cournot upper bound on
the markup.

\subsection{Marginal Revenue and the SFE First-Order Condition}
\label{app:foc_derivation}

Seller $i$'s profit at clearing price $p$ is
$\pi_i = (p - c) \cdot S_i(p)$, since marginal cost is constant at $c$.
Maximization is equivalent to maximizing profit along the residual demand
curve $R_i(\cdot)$ defined in equation~\eqref{eq:resid}. If seller $i$
considers increasing its offered quantity by $\mathrm{d}q$ at price $p$,
the market-clearing price must move by $\mathrm{d}p$ to preserve the
clearing identity~\eqref{eq:clearing}:
\begin{equation}
  \mathrm{d}p
  \;=\; \frac{\mathrm{d}q}{D'(p) - \sum_{j \neq i} S_j'(p) - S_f'(p)}.
  \label{eq:dprice}
\end{equation}
Because $D'(p) < 0$ and $S_j'(p) \geq 0$, the denominator is strictly
negative: an increase in supply lowers the clearing price. Marginal
revenue at $p$ is then
\begin{equation}
  \text{MR}_i(p)
  \;=\; p + \frac{S_i(p)}{D'(p) - \sum_{j \neq i} S_j'(p) - S_f'(p)}.
  \label{eq:mr}
\end{equation}
Setting $\text{MR}_i(p) = c$ and rearranging yields
equation~\eqref{eq:sfe_foc}, the system of $K$ coupled first-order ODEs
characterizing the supply function equilibrium. This is the $K$-seller
generalization of \citet{Green1992_british}'s equation~(3) for the
symmetric duopoly and of the general condition in
\citet{Anderson2008_sfe_asymmetric}.

\subsection{Symmetric Reduction}
\label{app:symmetric_derivation}

In a symmetric Nash equilibrium $S_i(p) = s(p)$ for all $i$. Substituting
into~\eqref{eq:sfe_foc} gives
\begin{equation}
  s(p) \;=\; -(p - c)
  \Bigl[D'(p) - (K-1)\,s'(p) - S_f'(p)\Bigr].
  \label{eq:sfe_sym_raw}
\end{equation}
Under Assumption~\ref{assm:fringe} the fringe is a step function, so
$S_f'(p) = 0$ almost everywhere. Solving~\eqref{eq:sfe_sym_raw} for
$s'(p)$ yields the symmetric SFE ODE~\eqref{eq:sfe_ode} reported in the
body, with boundary condition $s(\bar{p}) = \bar{q}$ from
Proposition~\ref{prop:holmberg}.

\subsection{Cournot Upper Bound}
\label{app:cournot}

The Cournot upper bound on the markup follows from the inverse-elasticity
rule applied at market demand:
\begin{equation}
  L^{\text{Cournot}} \;=\; \frac{1}{K \cdot |D'(p^*) \cdot p^* / D(p^*)|}.
  \label{eq:cournot_bound}
\end{equation}
This equals the residual-demand elasticity in the symmetric case and
provides a useful sanity check on the SFE-implied Lerner indices reported
in Section~\ref{sec:results}.
